\chapter{Analiza zahteva}

\section{Funkcionalni zahtevi}

Sistem za udomljavanje životinja mora zadovoljavati sledeće funkcionalne zahteve, grupisane po nivoima ocenjivanja:

\subsection{Nivo 6 – Osnovna funkcionalnost}

\begin{itemize}
    \item CRUD operacije za udruženja (svako udruženje ima administratora udruženja koji se dodaje prilikom kreiranja)
    \item CRUD operacije za životinje u okviru udruženja
    \item Prikaz svih udruženja (bez informacija o trošenju donacija)
\end{itemize}

\subsection{Nivo 7 – Autentifikacija i pregled}

\begin{itemize}
    \item Registracija volontera sa dodelom korisničkog imena i lozinke
    \item Prijava i odjava sa sistema za volontera i administratora udruženja
    \item Prikaz svih životinja iz nekog udruženja
\end{itemize}

\subsection{Nivo 8 – Upravljanje udomiteljima i privremeni smeštaj}

\begin{itemize}
    \item Uzimanje životinje u privremeni smeštaj (volonter)
    \item Prikaz svih udomitelja iz baze podataka (administrator udruženja i volonter)
    \item Dodavanje novog udomitelja u bazu podataka (administrator udruženja i volonter)
\end{itemize}

\subsection{Nivo 9 – Zahtevi za udomljavanje}

\begin{itemize}
    \item Slanje zahteva za udomljavanje životinje uz unos svojih kontakt podataka (neregistrovani korisnik)
    \item Prihvatanje ili odbijanje zahteva (administrator udruženja i volonter)
    \item Ako je zahtev prihvaćen, povezivanje životinje sa udomiteljem
\end{itemize}

\subsection{Nivo 10 – Napredne administrativne funkcije}

\begin{itemize}
    \item Pregled svih korisnika sistema (administrator)
    \item Pregled svih volontera iz udruženja (administrator udruženja)
    \item Mogućnost importa i pregleda izvoda računa udruženja (administrator udruženja)
\end{itemize}

\section{Nefunkcionalni zahtevi}

\begin{itemize}
    \item \textbf{Sigurnost}: Obavezne lozinke, različiti nivoi pristupa
    \item \textbf{Skalabilnost}: Sposobnost dodavanja novih udruženja i korisnika bez degradacije performansi
    \item \textbf{Održivost}: Jasna arhitektura i dokumentacija (UML dijagrami)
\end{itemize}
